\begin{frame}
\frametitle{Paměťová hierarchie od registrů k SSD}

{
\centering

\includegraphics[width=0.65\linewidth]{generated/memory/mem-reg-to-disc.pdf}

}
\vskip 2mm

\begin{tabular}{l|llll}
Typ      & L1 SRAM   & Sync SRAM &  DDR3      & HDD \\
Velikost & 32kB      & 1 MB      &  16 GB     & 3TB \\
Cena     & 10 kč/kB  & 300 kč/MB &  123 kč/GB & 1 kč/GB \\
Rychlost & 0.2...2ns & 0.5...8 ns/word & 15 GB/sec & 100 MB/sec \\
\end{tabular}

\vskip 2mm

Některá data mohou existovat ve více kopiích (úrovně, SMP).
Pro modifikaci dat je potřeba mechanizmů pro udržení koherence slov a konzistence datových struktur.

\end{frame}

\begin{frame}
\frametitle{Paměťová hierarchie -- základní principy}

\begin{itemize}
\item Programy/procesy přistupují v daném okamžiku jen k malé části svého adresového prostoru 
\item \textbf{Časová lokalita}
\begin{itemize}
\item Položky, ke kterým se přistupovalo nedávno, budou zapotřebí brzy znovu.
\item Příklad: programová smyčka, proměnné instrukcí.
\end{itemize}
\item \textbf{Prostorová lokalita}
\begin{itemize}
\item Položky poblíž právě používaným budou brzy zapotřebí také. 
\item Příklad: sekvenční přístup ke kódu (paměť programu), datová pole (paměť dat).
\end{itemize}
\end{itemize}
Princip se využívá jak v algoritmech (lokální proměnné), kompilátorech (přesun do registrů), na úrovni pamětí (automaticky), operačního systému (přesun z disku do paměti) případě opět programově, čtení a zápisy do souborů.
\end{frame}
